% ===========================================================================
% APPENDIX
% ===========================================================================

\section{Manual Forward-Pass Formulas for the \texttt{predict()} Bypass}
\label{sec:appendix}

This appendix reports the exact formulas used to compute each
surrogate's forward pass manually (Section~\ref{subsec:operational_feasibility}),
for reproducibility. In every case, weights are extracted once from
the trained model object into plain double-precision arrays, and the
forward pass is subsequently computed by direct matrix arithmetic with
no call to \texttt{predict()}, \texttt{minibatchpredict()}, or
construction of a \texttt{dlarray}/\texttt{RegressionGP} query object.

\subsection{Feedforward Networks (Residual MLP, SW-MLP, PINN-v2)}
\label{subsec:appendix_mlp}

For a network with $L$ fully-connected layers, weight matrices
$W_1,\ldots,W_L$, biases $b_1,\ldots,b_L$, and per-layer activations
$\phi_1,\ldots,\phi_L$ (identified automatically from the trained
\texttt{dlnetwork}'s layer graph: \texttt{tanhLayer},
\texttt{reluLayer}, or linear for the final layer), the forward pass on
normalized input $x_n$ is the standard recursion
\begin{equation}
h_0 = x_n, \qquad h_\ell = \phi_\ell\!\left(W_\ell h_{\ell-1} + b_\ell\right)
\text{ for } \ell=1,\ldots,L, \qquad y = h_L.
\label{eq:appendix_mlp}
\end{equation}
SW-MLP additionally flattens its $[\text{seq\_len}\times 4]$ sliding
window into a single input vector via column-major reshaping before
Eq.~\eqref{eq:appendix_mlp} is applied, matching the ordering used at
training time. Agreement with \texttt{predict()} was verified to
$1.4\times10^{-7}$ (SW-MLP), $2.2\times10^{-7}$ (PINN-v2), and
$1.0\times10^{-7}$ (the residual MLP) across multiple random test
points in the training domain.

\subsection{Gaussian Process (GP-v2)}
\label{subsec:appendix_gp}

With a Mat\'ern~5/2 kernel, constant basis function, active set
$\{\mathbf{x}_i\}_{i=1}^{N}$, dual coefficients $\{\alpha_i\}$, and
constant basis coefficient $\beta_0$ (all extracted directly from the
trained \texttt{RegressionGP} object's \texttt{ActiveSetVectors},
\texttt{Alpha}, and \texttt{Beta} properties), the posterior mean at a
normalized query point $\mathbf{x}$ is
\begin{equation}
m(\mathbf{x}) = \beta_0 + \sum_{i=1}^{N} \alpha_i\, k(\mathbf{x}, \mathbf{x}_i),
\qquad
k(\mathbf{x},\mathbf{x}') = \sigma_f^2\!\left(1+\frac{\sqrt{5}\,r}{l}+\frac{5r^2}{3l^2}\right)
e^{-\sqrt{5}r/l},
\label{eq:appendix_gp}
\end{equation}
with $r=\lVert\mathbf{x}-\mathbf{x}'\rVert_2$ and kernel hyperparameters
$(\sigma_f, l)$ read from \texttt{KernelInformation.KernelParameters}.
Reducing the active-set size from $N_{\text{sub}}=400$ (used for the
open-loop accuracy results of Section~\ref{subsubsec:gp_v2}) to
$N_{\text{sub}}=50$ for the closed-loop diagnostic
(Section~\ref{subsec:operational_feasibility}) reduced manual per-step
cost roughly in proportion to $N_{\text{sub}}$, with no measurable loss
of test-set accuracy on this task. Agreement with \texttt{predict()}
was verified to $4.2\times10^{-10}$.

\subsection{Temporal Convolutional Network (TCN)}
\label{subsec:appendix_tcn}

For each of the three dilated causal convolution blocks (kernel size
$K=3$, dilations $1,2,4$, input channels $C_{\text{in}}$, output
channels/filters $C_{\text{out}}=16$), the layer's weight tensor
$W\in\mathbb{R}^{K\times C_{\text{in}}\times C_{\text{out}}}$ is
reshaped once, at extraction time, into
$W_{\text{mat}}\in\mathbb{R}^{C_{\text{out}}\times KC_{\text{in}}}$ by
horizontally concatenating $K$ blocks $W_{\text{mat}}^{(k)} =
W(k,:,:)^\top$. At inference, the causally left-padded input
$X_p\in\mathbb{R}^{C_{\text{in}}\times(L+d(K-1))}$ (zero-padded by
$d(K-1)$ steps, $d$ the block's dilation) is unfolded into
$U\in\mathbb{R}^{KC_{\text{in}}\times L}$ by stacking $K$ dilated
slices, and the entire convolution is computed as the single matrix
product
\begin{equation}
Y = W_{\text{mat}}\, U + b,
\label{eq:appendix_tcn}
\end{equation}
followed by per-timestep, channel-wise layer normalization
($Y_{:,t} \leftarrow \gamma\odot(Y_{:,t}-\mu_t)/\sqrt{\sigma_t^2+\epsilon}+\delta$,
with $\mu_t,\sigma_t^2$ computed across channels at each timestep $t$)
and ReLU. An initial implementation of Eq.~\eqref{eq:appendix_tcn} as
explicit triple-nested loops over filters, timesteps, and kernel taps
was numerically correct but yielded only a $2.2\times$ speedup over
\texttt{predict()}; precomputing $W_{\text{mat}}$ once and using
Eq.~\eqref{eq:appendix_tcn} as a single matrix multiplication per block
was necessary to reach the $56\times$ speedup reported in
Table~\ref{tab:predict_bypass}. Agreement with \texttt{predict()} was
verified to $2.2\times10^{-7}$ across five random test points.

\subsection{LSTM}
\label{subsec:appendix_lstm}

For a single LSTM layer with hidden size $H$, input $x_t$, and
previous hidden/cell state $(h_{t-1}, c_{t-1})$, MATLAB's
\texttt{lstmLayer} concatenates gate weights in the order input,
forget, cell-candidate, output (each occupying $H$ consecutive rows of
the $4H$-row \texttt{InputWeights} $W_i$, \texttt{RecurrentWeights}
$W_r$, and \texttt{Bias} $b$). One recurrent step is
\begin{align}
z &= W_i x_t + W_r h_{t-1} + b, \qquad z\in\mathbb{R}^{4H}, \label{eq:appendix_lstm1}\\
i_t &= \sigma(z_{1:H}), \quad f_t = \sigma(z_{H+1:2H}), \quad
g_t = \tanh(z_{2H+1:3H}), \quad o_t = \sigma(z_{3H+1:4H}), \\
c_t &= f_t \odot c_{t-1} + i_t \odot g_t, \qquad
h_t = o_t \odot \tanh(c_t), \label{eq:appendix_lstm3}
\end{align}
with $\sigma(\cdot)$ the logistic sigmoid. The trained two-layer
network (hidden sizes $64$ and $32$) applies
Eqs.~\eqref{eq:appendix_lstm1}--\eqref{eq:appendix_lstm3} recurrently
over the $L=10$-step input window for layer 1 (retaining every $h_t$ as
input to layer 2), then again over the resulting sequence for layer 2
(retaining only its final $h_L$), followed by a linear output
projection. Both hidden and cell states are initialized to zero at the
start of each window, matching \texttt{predict()}'s default (stateless)
behavior. Unlike the other architectures, this recursion cannot be
reduced to a single matrix multiplication, since $h_t$ depends on
$h_{t-1}$; the ten sequential steps are an irreducible cost of the
manual implementation, consistent with LSTM's manual per-step cost
($\SI{52.2}{ms}$) remaining $2$--$4\times$ that of the five single-pass
architectures even after the bypass. Agreement with \texttt{predict()}
was verified to $1.3\times10^{-5}$--$2.0\times10^{-4}$ across five
random test points  an order of magnitude looser than the
$10^{-7}$--$10^{-13}$ agreement obtained for the feedforward
architectures and GP, attributed to accumulated single-precision
(\texttt{dlnetwork}'s default) versus double-precision (manual
implementation) rounding compounding over ten sequential nonlinear
steps, rather than to an incorrect gate ordering, since the
discrepancy is small and consistently sized across independent test
points rather than an order-of-magnitude mismatch.

\subsection{Reproducibility: Explicit Random Seeding}
\label{subsec:appendix_seeding}

All V2 surrogate training scripts (LSTM, TCN, PINN-v2, SW-MLP) fix
MATLAB's global random state via \texttt{rng(cfg.data.split\_seed)}
immediately before network construction and training, and record the
seed used in the returned model struct (\texttt{mdl.rng\_seed}). This
addresses an earlier version of this pipeline, in which weight
initialization and mini-batch shuffling were left to MATLAB's
unseeded global random state, so that repeated training runs produced
qualitatively similar but numerically different results. The Table~\ref{tab:stage2_accuracy}
V2 figures were regenerated under this explicit seeding and confirmed
reproducible across repeated runs on the same MATLAB release and
machine; exact reproducibility across different releases or hardware
is not guaranteed, since \texttt{trainnet}'s underlying numerical
kernels may differ across releases.

\subsection{Consolidated Verification Tolerances}
\label{subsec:appendix_verification_table}

The manual-bypass forward pass for each architecture was verified
against \texttt{predict()} independently, at different points in this
diagnostic (Sections~\ref{subsec:appendix_mlp}--\ref{subsec:appendix_lstm}
above, and the mechanism-isolation test of
Section~\ref{subsec:operational_feasibility}). Where an architecture
was verified more than once (TCN: $2.05\times10^{-7}$ and
$2.18\times10^{-7}$ in two independent runs), Table~\ref{tab:verification_tolerances}
reports the single largest (worst-case) deviation observed, adopting
\emph{maximum absolute deviation} as the sole reporting convention
throughout this paper.

\begin{table}[H]
\centering
\caption{Consolidated manual-vs-\texttt{predict()} verification
tolerances, maximum absolute deviation, one entry per architecture.}
\label{tab:verification_tolerances}
\begin{tabular}{lccc}
\toprule
Architecture & Max.\ abs.\ deviation & Test points & Precision \\
\midrule
MLP-residual & $1.0\times10^{-7}$ & 1 & double \\
SW-MLP       & $1.4\times10^{-7}$ & 1 & double \\
PINN-v2      & $2.2\times10^{-7}$ & 1 & double \\
GP-v2        & $4.2\times10^{-10}$ & 1 & double \\
TCN          & $2.2\times10^{-7}$ & 5 & double \\
LSTM         & $2.0\times10^{-4}$ & 5 & double (vs.\ network's
  native single) \\
\bottomrule
\end{tabular}

\vspace{4pt}
{\footnotesize LSTM's tolerance is an order of magnitude looser than
the other five architectures, attributed to single- versus
double-precision accumulation over ten sequential recurrent steps
(Section~\ref{subsec:appendix_lstm}), not a logical error.}
\end{table}

\subsection{Diagnostic Note on Solver \texttt{ExitFlag}}
\label{subsec:appendix_exitflag}

Practitioners replicating this diagnostic should note that
\texttt{nlmpcmove}'s \texttt{info.ExitFlag} being non-positive (solver
iteration limit reached without formal convergence) is, under the
bounded-iteration solver configuration used throughout this study
(\texttt{MaxIterations}$=30$, Section~\ref{subsec:nlmpc_formulation}),
expected for \emph{every} controller at \emph{every} step, including
the fast and well-behaved Baseline. It is therefore not, by itself, a
useful signal of surrogate-specific infeasibility, and should not be
interpreted as such without also examining per-step CPU cost and the
resulting closed-loop trajectory.
